\documentclass[a4paper,11pt]{article}

\pagestyle{empty}
\setlength{\parindent}{0cm}

% Pacotes Principais -----------------------------------------------------------
\usepackage[portuges,brazil]{babel}
\usepackage[utf8]{inputenc}

% Figuras e Imagens ------------------------------------------------------------
\usepackage{graphicx}
% Figuras lado a lado
\usepackage{epsfig}
\usepackage{subfigure}

% Utilizar H para inserir as imagens REALMENTE onde eu desejo
\usepackage{float}

% Fontes -----------------------------------------------------------------------
\usepackage[T1]{fontenc}
\usepackage{pslatex}

% Simbolos ---------------------------------------------------------------------
\usepackage{textcomp}

% Tabelas ----------------------------------------------------------------------
%\usepackage{multicol}
\usepackage{multirow}
% Colorir a tabela
\usepackage{colortbl}

% Outros pacotes ---------------------------------------------------------------
\usepackage{noitemsep}

\usepackage{color}
\usepackage{xcolor}

% Comentários em bloco
\usepackage{verbatim}

% Sublinhado, traçado...
\usepackage{ulem}

% Desenhar
\usepackage{tikz}

% Comandos gerais --------------------------------------------------------------

% Configuração da fonte
\renewcommand{\familydefault}{\sfdefault}

% Comandos matemáticos ---------------------------------------------------------
% Implicação em fórmulas
\newcommand{\implica}{$\quad\Rightarrow\quad$} %Meio de linha
\newcommand{\implicafim}{$\quad\Rightarrow$}   %Fim de linha
\newcommand{\tende}{$\rightarrow$}

% Fração com parenteses
\newcommand{\pfrac}[2]{\parent{\frac{#1}{#2}}}

% Transformada de Laplace e transformada Z
\newcommand{\lapl}{\pounds}
\newcommand{\transfz}{\mathcal{Z}}

% Sequências
\newcommand{\sequencia}[4]{$#1_{#2}$, $#1_{#3}$, \ldots, $#1_{#4}$}

% Outros ----------------------------------------------------------------------
\newcommand{\chave}[1]{\left\{#1\right\}}
\newcommand{\colchete}[1]{\left[#1\right]}
\newcommand{\parent}[1]{\left(#1\right)}

\let\D\displaymath

\newtheorem{definicao}{Definição}
\newtheorem{exemplo}{Exemplo}
\newtheorem{lema}{Lema}
\newtheorem{observ}{Observação}
\newtheorem{teorema}{Teorema}

\newcommand{\defin}[1]{\begin{definicao}#1\end{definicao}}
\newcommand{\exemp}[1]{\begin{exemplo}#1\end{exemplo}}
\newcommand{\obs}[1]{\begin{observ}#1\end{observ}}

\newcommand{\pixel}{{\it pixel}}

% Tikz
\usetikzlibrary{matrix}


\begin{document}

\begin{table}[H]
\centering
\begin{tabular}{cl}
% Primeira Linha
\multirow{7}{*}
{
\includegraphics[width=1.65cm]{imgs/uern}
} & \\
& Governo do Estado do Rio Grande do Norte\\
% Segunda Linha
& Secretaria de Estado da Educação, da Cultural e dos Desportos -- SECD\\
% Terceira Linha
& {\sc Universidade do Estado do Rio Grande do Norte -- UERN}\\
% Quarta Linha
& Pró-Reitoria de Ensino e Graduação -- PROEG\\
% Quinta Linha
& Ciências da Computação\\
& 
\end{tabular}
\end{table}

\begin{center}
\uline{{\sc Lista de Exercícios}}
\end{center}

\vspace{0.5cm}

1) Defina o que são mídias de percepção e mídias de representação.

\vspace{0.25cm}

2) Sabendo que a informação é representada a partir de variações das
características dos sinais, determine que tipos de problemas podem ocorrer na
transmissão de um sinal.

\vspace{0.25cm}

3) Diferencie a representação de imagens no formato vetorial e no formato
matricial identificando suas características.

\vspace{0.25cm}

4) Descreva passo a passo, em detalhes, como se dá a conversão A/D para um sinal
analógico de 11 KHz, utilizando 1 Byte para representar cada amostra. Explique
também porque ocorre o erro de quantização.

\vspace{0.25cm}

5) O que é resolução geométrica? E resolução de cor?

\vspace{0.25cm}

6) Qual será o tamanho do arquivo (em KBytes) de uma imagem de $30 \times 30$
{\it pixels} que utiliza o modelo de cores RGB (8 bits/cor)?

\vspace{0.25cm}

7) Explique o que é multimídia e identifique quais são os recursos
computacionais que podem ser utilizados por aplicativos multimídia.

\vspace{0.25cm}

8) Descreve quais são as principais características dos arquivos multimídia.

\vspace{0.25cm}

9) O que são títulos, aplicativos e sítios? Diferencie títulos lineares de
títulos hipermídia.

\vspace{0.25cm}

10) Determine e descreva quais são os tipos de plataformas.

\vspace{0.25cm}

11) Identifique as características das famílias de plataformas Windows, Linux e
MacOS.

\vspace{0.25cm}

12) Qual é o objetivo da utilização de bibliotecas e qual é o princípio de
funcionamento?

\vspace{0.25cm}

13) O que significa a sigla OLE e quais são suas características?

\vspace{0.25cm}

14) Por que as aplicações fechadas possuem, normalmente, uma curva de
aprendizado mais rápida?

\vspace{0.25cm}

15) As linguagens de programação permitiram maior flexibilidade para a criação
de aplicações multimídia. Descreva algumas características das linguagens C,
C++, Java e Python.

\vspace{0.25cm}

16) Explique o que são e dê exemplos de ferramentas híbridas.

\vspace{0.25cm}

17) Descreva as características da navegação não-linear.

\vspace{0.25cm}

18) O que são arquivos HTML, hiperdocumentos locais, padrão SMIL, sítios
estáticos e sítios dinâmicos?

\vspace{0.25cm}

19) Identifique cada um dos elementos presentes em projetos multimídia.

\vspace{0.25cm}

20) Identifique e descreva as atividades realizadas durante a fase de
desenvolvimento no modelo de ciclo de vida.

\vspace{0.25cm}

21) Por que motivo as equipes de desenvolvimento, normalmente, são compostas por
tantas pessoas? Identifique o papel de cada uma dessas.

\vspace{0.25cm}

22) Dê exemplos de algumas restrições que devem ser observadas que estão
relacionadas com os direitos autorais.

\vspace{0.25cm}

23) No processo técnico a especificação compreende duas fases. Determine quais
são essas fases e cite suas características.

\vspace{0.25cm}

24) Sobre o que trata o desenho externo na fase de desenvolvimento? E o desenho
de arquitetura?

\vspace{0.25cm}

25) Na representação digital de imagens, o que é resolução espacial?

\vspace{0.25cm}

26) Quais são os elementos que compõem a visão humana?

\vspace{0.25cm}

27) Como funcionam os monitores CRT convencionais?

\vspace{0.25cm}

28) Dê exemplos de dispositivos de saída do tipo ``cópia permanente''.

\vspace{0.25cm}

29) No processamento digital de imagens existem, basicamente, dois tipos de
domínio de operação. Identifique-os e cite exemplos de processamentos em cada um
deles.

\vspace{0.25cm}

30) Diferencie a compressão de imagens com perda da compressão sem perda e cite
alguns dos algoritmos e formatos de imagem mais conhecidos.

\end{document}
